\notechapter{N-20221218}{Basic Topology I: Metric Spaces, Boundedness, Norms, Open Sets, and Completeness}

\section{Why the note starts with metrics}

This note begins a topology sequence, but it does not begin with arbitrary open sets.  It begins with metric spaces, because metric spaces keep the analysis intuition alive.  In real analysis, convergence means
\[
  x_n\to x
  \quad\Longleftrightarrow\quad
  (\forall \varepsilon>0)(\exists N)(\forall n>N)\ |x_n-x|<\varepsilon.
\]
The first abstraction is to replace $|x_n-x|$ by a general distance $d(x_n,x)$.  Thus a sequence in a metric space converges when
\[
  (\forall \varepsilon>0)(\exists N)(\forall n>N)\ d(x_n,x)<\varepsilon.
\]
The hidden question is: which parts of analysis really require a numerical distance, and which only require a notion of neighborhood?  Metric spaces are the intermediate language between analysis and topology.

\section{Boundedness: diameter and radius}

A metric space $M$ is bounded if there is a constant $D$ such that
\[
  d(p,q)\le D\qquad \text{for all }p,q\in M.
\]
One can think of $D$ as a diameter bound.  The note also gives an equivalent radius formulation: fixing a point $p$, if there is $R$ such that
\[
  d(p,q)\le R\qquad \text{for all }q\in M,
\]
then $M$ is bounded.  Indeed, for any $q_1,q_2\in M$,
\[
  d(q_1,q_2)\le d(q_1,p)+d(p,q_2)\le 2R.
\]
Conversely, a diameter bound immediately gives a radius bound from any chosen point.

The examples are deliberately chosen to break naive intuition: $[0,1]$ is bounded, $\mathbb R^n$ is not, an infinite set with the discrete metric is bounded, but $\mathbb N$ with the usual metric is not.  Thus boundedness belongs to a metric, not to a bare set.

\section{Metric spaces}

A metric space is a pair $(X,d)$ where $d:X\times X\to \mathbb R$ satisfies:
\begin{enumerate}[(i)]
  \item $d(x,y)\ge0$;
  \item $d(x,y)=0$ iff $x=y$;
  \item $d(x,y)=d(y,x)$;
  \item $d(x,z)\le d(x,y)+d(y,z)$.
\end{enumerate}
A subset $Y\subseteq X$ inherits the subspace metric
\[
  d_Y(a,b)=d_X(a,b).
\]
This is the metric ancestor of subspace topology.

\section{Sequences and continuity}

A sequence $(x_n)$ converges to $x$ if every small ball around $x$ eventually contains all $x_n$.  In $\mathbb R^k$ with the Euclidean metric, convergence is coordinatewise:
\[
  x_n=(x_n^1,\ldots,x_n^k)\to x=(x^1,\ldots,x^k)
\]
iff $x_n^i\to x^i$ for each $i$.  One direction follows from $|x_n^i-x^i|\le \|x_n-x\|_2$; the other follows by summing the small coordinate errors.

A map $f:X\to Y$ between metric spaces is continuous at $x_0$ if
\[
  (\forall \varepsilon>0)(\exists \delta>0)\quad
  d_X(x,x_0)<\delta \Rightarrow d_Y(f(x),f(x_0))<\varepsilon.
\]
In metric spaces, continuity can also be detected by sequences:
\[
  x_n\to x_0 \Rightarrow f(x_n)\to f(x_0).
\]
The later general topology notes warn that sequences no longer detect everything in arbitrary topological spaces.

\section{Standard metrics}

On $\mathbb R^n$ the Euclidean, taxicab, and max metrics are
\[
  d_2(x,y)=\left(\sum_i |x_i-y_i|^2\right)^{1/2},
\]
\[
  d_1(x,y)=\sum_i |x_i-y_i|,
\]
and
\[
  d_\infty(x,y)=\max_i |x_i-y_i|.
\]
Their balls have different shapes: round disks, diamonds, and squares.  Nevertheless they generate the same usual topology on $\mathbb R^n$.  This is an early hint that topology remembers neighborhoods but forgets exact metric geometry.

The discrete metric on a set $X$ is
\[
  d(x,y)=\begin{cases}
  0,&x=y,\\
  1,&x\ne y.
  \end{cases}
\]
Then $B_{1/2}(x)=\{x\}$.  Every subset is open and closed, and a sequence converges iff it is eventually constant.

\section{Norms and inner products}

A norm on a vector space $V$ satisfies positivity, homogeneity, and the triangle inequality:
\[
  \|v+w\|\le \|v\|+\|w\|.
\]
Every norm gives a metric by
\[
  d(v,w)=\|v-w\|.
\]
The $p$-norms are
\[
  \|x\|_p=\left(\sum_i |x_i|^p\right)^{1/p},
\]
and $\|x\|_\infty=\max_i |x_i|$.

An inner product gives a norm by
\[
  \|v\|=\sqrt{\langle v,v\rangle}.
\]
The triangle inequality for this norm follows from Cauchy--Schwarz.  So Euclidean metric geometry, linear algebra, and inequalities are already intertwined.

\section{Open and closed balls}

For $x\in X$ and $r>0$,
\[
  B_r(x)=\{y:d(x,y)<r\},\qquad
  \overline B_r(x)=\{y:d(x,y)\le r\}.
\]
Open balls are open.  If $y\in B_r(x)$, choose
\[
  \delta=r-d(x,y)>0.
\]
Then $B_\delta(y)\subseteq B_r(x)$ by the triangle inequality.  The proof says: every point inside a ball has some remaining margin before hitting the boundary.

\section{Open sets, closed sets, and limit points}

A set $U\subseteq X$ is open if every point $x\in U$ has some ball $B_r(x)$ contained in $U$.  A set is closed if its complement is open.  In metric spaces, a set $C$ is closed iff it contains the limits of all convergent sequences from $C$.

A point $x$ is a limit point of $A$ if every ball around $x$ meets $A$ in a point other than $x$.  The closure satisfies
\[
  x\in \overline A
  \quad\Longleftrightarrow\quad
  (\forall r>0)\ B_r(x)\cap A\ne\varnothing.
\]
The boundary is
\[
  \partial A=\overline A\cap \overline{X\setminus A}.
\]
Examples matter: every real number is a limit point of $\mathbb Q$, and also of the irrational numbers; in the discrete metric, points are isolated by small balls.

\section{Cauchy sequences and completeness}

A sequence is Cauchy if
\[
  (\forall \varepsilon>0)(\exists N)(\forall m,n>N)\ d(x_m,x_n)<\varepsilon.
\]
A metric space is complete if every Cauchy sequence converges in the space.  The standard contrast is
\[
  \mathbb R \text{ is complete},\qquad \mathbb Q \text{ is not complete}.
\]
A rational sequence converging to $\sqrt2$ is Cauchy in $\mathbb Q$, but it has no rational limit.

Convergence names a destination; Cauchy only says the sequence internally stabilizes.  Completeness says internal stabilization is enough to guarantee a destination inside the space.

\section{Total boundedness and compactness}

The note also points toward total boundedness and compactness.  A metric space is totally bounded if for every $\varepsilon>0$, finitely many $\varepsilon$-balls cover it.  Total boundedness is stronger than boundedness: an infinite discrete metric space is bounded, but for $\varepsilon<1$ it cannot be covered by finitely many $\varepsilon$-balls.

In $\mathbb R^n$, compactness is equivalent to closed and bounded, but that is the Heine--Borel theorem, not the definition in arbitrary metric spaces.  A better metric-space memory is:
\[
  \text{compact} \Rightarrow \text{complete and totally bounded}.
\]
Conversely, in metric spaces, complete plus totally bounded implies compact.

\section{Beyond the first calculation}

This note stages the transition from analysis to topology.  Metrics let us talk about boundedness, convergence, continuity, balls, closed sets, limit points, Cauchy sequences, and completeness.  But later topology keeps only what can be expressed through open sets.  Boundedness and completeness will fail to be topological invariants; continuity and connectedness will survive.


\paragraph{Expanded synthesis.}
The common theme of this chapter is that an elementary limiting or computational rule becomes powerful only after one specifies the space in which the limiting process takes place. The earlier sections (`Why the note starts with metrics`, `Boundedness: diameter and radius`, `Metric spaces`, `Sequences and continuity`, `Standard metrics`) may look like separate techniques, but they all ask the same question: which operations are stable under approximation? A derivative is stable first-order approximation,
\[
  f(a+h)=f(a)+L(h)+o(|h|),
\]
while an integral is stable accumulation with respect to a measure, and a convergent series is a limit in a chosen norm or topology.

The advanced bridge is functional-analytic. Uniform convergence is convergence in a sup norm; $L^p$ convergence is convergence in an integral norm; differentiability is the existence of a continuous linear approximation; many differential equations are operator equations $L[u]=f$. Theorems such as dominated convergence, the inverse function theorem, the projection theorem in Hilbert spaces, and semigroup existence results are refined answers to the same elementary concern: when may we pass from finite approximations to a genuine limiting object?

To use this viewpoint when rereading the original examples, identify the hidden stability condition. A Taylor expansion asks for control of the remainder. A change of variables asks for differentiability and Jacobian control. Termwise differentiation of a series asks for uniform convergence of derivative series. A differential-equation formula asks whether the operator is invertible under the imposed initial or boundary data. The elementary computation is the visible part; the advanced theory names the hypotheses that make it legitimate.
